\section{Models and Experimental Setup}
\label{sec:models_setup}

\subsection{Model Architectures}
We train three distinct supervised learning algorithms, each fitted independently \emph{once per endpoint} exclusively on that endpoint's labeled training rows. A molecule with an unmeasured label for an endpoint is strictly omitted from both model fitting and performance evaluation for that specific task; missing labels are never imputed as inactive negatives. Every model is evaluated on two feature spaces: \texttt{morgan} (2,048-bit, radius-2 circular fingerprints) and \texttt{morgan+desc} (the 2,048 Morgan bits concatenated with 213 preprocessed, scaled RDKit descriptors; Section~\ref{sec:features}).

\begin{itemize}
    \item \textbf{Logistic Regression (LR):} Implemented via Scikit-learn's~\cite{pedregosa2011scikit} \texttt{liblinear} coordinate descent solver with an L1 or L2 regularization penalty and a maximum iteration threshold of 2,000. Because liblinear with fixed tolerance is completely deterministic, seed variance across model fits is exactly $0.0000$.
    \item \textbf{Random Forest (RF):} Implemented using Breiman's~\cite{breiman2001random} ensemble of decorrelated classification trees via Scikit-learn's \texttt{RandomForestClassifier}. Splitting criteria evaluate Gini impurity across $\sqrt{d}$ candidate features per split.
    \item \textbf{XGBoost (XGB):} Implemented via Chen and Guestrin's~\cite{chen2016xgboost} extreme gradient boosting system using histogram-optimized tree building. To prevent overfitting on sparse features, row subsampling is fixed at $0.80$, column subsampling per tree is set to $0.50$, and PR AUC (\texttt{aucpr}) is designated as the internal objective. Tree counts are incorporated into the tuning grid without early stopping to ensure strict separation of validation tuning from test evaluation.
\end{itemize}

\subsection{Class Imbalance Mitigation}
To address severe positive sparsity without distorting underlying molecular feature distributions through synthetic sampling (e.g., SMOTE), we employ adaptive loss weighting computed from each assay's observed training labels ($n_{\mathrm{pos}}$ and $n_{\mathrm{neg}}$):
\begin{itemize}
    \item LR applies inverse frequency weighting: $\texttt{class\_weight="balanced"}$.
    \item RF applies dynamic bootstrap weighting: $\texttt{class\_weight="balanced\_subsample"}$, recomputing class weights on each bootstrap sample.
    \item XGB scales positive gradient steps: $\texttt{scale\_pos\_weight} = n_{\mathrm{neg}} / n_{\mathrm{pos}}$.
\end{itemize}

\subsection{Tuning Protocol}
To maintain rigorous fairness, every model family is allocated an identical search budget of exactly eight hyperparameter configurations (Table~\ref{tab:grids}). All configurations are trained on the scaffold training partition ($n=6{,}167$) and evaluated on the scaffold validation split ($n=771$) with seed 0. For each model family and feature modality, the configuration yielding the highest mean validation PR AUC across the 12 endpoints was selected (Table~\ref{tab:chosen}). 

The selected hyperparameter configurations were serialized to \texttt{results/best\_configs.json} and immutably committed to git version control (commit \texttt{041892c}) prior to executing any evaluation against the test partition. In accordance with the project contract, configurations tuned on the scaffold split were applied unchanged to the random split, guaranteeing that the test partition was touched strictly once during final evaluation.

\begin{table}[t]
\caption{Pre-registered hyperparameter tuning grids (8 configurations per model).}
\label{tab:grids}
\centering\small
\begin{tabular}{@{}lp{0.78\columnwidth}@{}}
\toprule
Model & Configurations \\
\midrule
LR  & $C \in \{0.01, 0.03, 0.1, 0.3, 1, 3\}$ with L2; $C \in \{0.1, 1\}$ with L1 \\
RF  & 300 trees $\times$ max depth $\in\{\text{none}, 30\}$ $\times$ min leaf $\in\{1, 3, 5\}$; plus (500 trees, none, 3) and (500, 30, 5) \\
XGB & (max depth, learning rate, trees): (3, 0.10, 300), (5, 0.10, 300), (7, 0.10, 300), (3, 0.05, 500), (5, 0.05, 500), (7, 0.05, 500), (4, 0.10, 400), (6, 0.03, 600) \\
\bottomrule
\end{tabular}
\end{table}

\begin{table}[t]
\caption{Selected hyperparameter configurations and mean validation PR AUC on scaffold split (seed 0). RF: (trees, max depth, min leaf); XGB: (max depth, learning rate, trees).}
\label{tab:chosen}
\centering\small
\setlength{\tabcolsep}{4pt}
\begin{tabular}{@{}llll@{}}
\toprule
Model & Feature Set & Selected Hyperparameters & Val.\ PR AUC \\
\midrule
LR  & morgan      & L2 penalty, $C=0.03$       & 0.3475 \\
LR  & morgan+desc & L2 penalty, $C=0.03$       & 0.3429 \\
RF  & morgan      & (300, none, 1)             & 0.4000 \\
RF  & morgan+desc & (300, none, 1)             & 0.4433 \\
XGB & morgan      & (6, 0.03, 600)             & 0.3607 \\
XGB & morgan+desc & (7, 0.05, 500)             & 0.4307 \\
\bottomrule
\end{tabular}
\end{table}

\subsection{Evaluation Methodology and Metric Masking}
For each endpoint, model predictions are evaluated using Receiver Operating Characteristic Area Under the Curve (ROC AUC) and Precision--Recall Area Under the Curve (PR AUC), calculated strictly over labeled test rows. If an evaluation partition possesses fewer than two distinct classes for a target, the metric is undefined (\texttt{NaN}) and omitted from aggregate averaging.

Because active compounds are rare, PR AUC is reported as the primary metric of clinical and chemical relevance: while ROC AUC can appear inflated when true negatives dominate, PR AUC directly reflects the positive predictive value as a function of recall.

Final test experiments are executed across seeds 0, 1, and 2 on both the scaffold split and the random split. We report the mean and sample standard deviation across the three seeds. The full experimental design comprises:
\begin{equation*}
    3\text{ models} \times 2\text{ feature sets} \times 2\text{ splits} \times 3\text{ seeds} \times 12\text{ endpoints} = 432
\end{equation*}
fitted classifiers, all orchestrated through a unified, leak-free evaluation harness.
